\section{Conclusion}

This paper frames MRW-like inference from finite stochastic signals as a
validity-aware spectral diagnostic problem. CMIN-SR separates empirical
spectrum estimation from spectral interpretation, compares monofractal and MRW
projection families, and reports diagnostic scores for scaling, curvature,
monofractal compatibility, MRW compatibility, boundary behavior, and
instability.

The experiments support a balanced conclusion. In analytic spectrum space, the
MRW and monofractal geometries are separable and stable across random seeds.
In raw finite samples, however, empirical spectrum estimation is the limiting
step: deterministic and neural zeta estimators do not reliably preserve the
MRW curvature coordinate under the tested short-window settings. The added
classical baseline comparison strengthens this interpretation: MFDFA improves
some metrics relative to simple structure functions, but MFDFA, aggregate
structure functions, and lightweight Haar wavelet controls do not produce
stable short-window \(\lambda^2\) recovery on the tested grid. The real-data
sanity check further shows that the framework can be applied to complex
financial factor returns, but only as a diagnostic tool, not as a mechanism
test.

CMIN-SR should therefore be used to organize evidence for or against
MRW-compatible spectral geometry. It should not be used as a guaranteed
short-window \(\lambda^2\) recovery engine or as proof that an observed signal
was generated by an MRW process.
